\section{Introduction}
Large Language Models~(LLMs) have evolved from specialized text processing tools into a widely applied technology in various fields\nobreak\cite{LLM_survey1,LLM_survey2,LLM_survey3,LLM_survey4,LLM_survey5}.
They have been applied in fields such as natural language processing\nobreak\cite{NLP_survey1,NLP_survey2,NLP_survey3}, code generation\nobreak\cite{code_generation1,code_generation2,code_generation3}, and conversational AI\nobreak\cite{conversational_agents1,conversational_agents2}.
The development of multimodal LLMs, which integrate modalities such as text, images, and audio, is further expanding their scope of application on a large scale\nobreak\cite{multi-modal1,multi-modal4,multi-modal5}. 
The interaction between algorithmic development and computational architectures enables these innovations.
Therefore, improving this infrastructure is not merely a practical challenge but critical for the broader societal integration of AI.
Currently, this technological progress primarily depends on the evolution of General-Purpose Graphics Processing Units~(GPGPUs).
However, despite these advances, achieving high performance with LLMs requires substantial computational resources for both training and inference.
In particular, the high performance of GPGPUs requires substantial power, which creates significant environmental and economic concerns\nobreak\cite{gpu_power1,gpu_power2}. 
The International Energy Agency~(IEA) estimates that data center electricity consumption could double by 2030, reaching approximately~\SI{945}{\TWh}\nobreak\cite{iea_energy_ai}. 
This consumption level, fueled primarily by the expanding demand for AI, slightly exceeds Japan's total annual electricity consumption.
This sustainability challenge presents a significant barrier to the widespread adoption of LLM technology, making it necessary to address this with improvements in computational architecture.

To address this challenge, specialized hardware such as Application-Specific Integrated Circuits~(ASICs) and Field-Programmable Gate Arrays~(FPGAs) have been widely investigated.
While edge GPUs, such as the NVIDIA Jetson series, demonstrate progress in reducing power consumption, their general-purpose graphics pipelines inherently constrain substantial gains in power efficiency.
In contrast, ASICs can potentially achieve orders of magnitude higher power efficiency. 
ASICs specialize their circuits for the core computational patterns of LLM inference, such as dot-product operations. 
This allows them to remove all unnecessary functionality and achieve maximum performance for a specific task.
However, this high efficiency is intrinsically coupled with a lack of flexibility, rendering these designs unable to adapt to evolving algorithms.


Therefore, we consider Coarse-Grained Reconfigurable Arrays~(CGRAs) as an architectural paradigm that addresses the fundamental trade-off between efficiency and flexibility.
CGRAs achieve power efficiency approaching that of an ASIC by mapping dataflow graphs directly onto an array of processing elements~(PEs), while preserving programmability through their reconfigurable fabric.
This work focuses on IMAX\nobreak\cite{imax3}, a general-purpose accelerator based on Coarse-Grained Linear Arrays~(CGLAs), a type of CGRA architecture.
IMAX features a linear arrangement of PEs and Local Memory Modules~(LMMs), a design specifically intended to efficiently handle the irregular memory access patterns of LLM inference. 


Our previous work has demonstrated the versatility and effectiveness of the IMAX architecture across a diverse range of workloads, including Sparse General Matrix Multiplication~(SpGEMM), Fast Fourier Transform~(FFT), and Convolutional Neural Networks~(CNNs)\nobreak\cite{imax3,cnn_imax_1,cnn_imax_2}. 
We have also extended this versatility to LLMs, showing the feasibility of implementing key kernels from the Llama2-based model on IMAX\nobreak\cite{first_llm_imax,llama2_imax}. 
Based on this foundation, this work provides the first comprehensive evaluation of IMAX as an LLM accelerator.
Specifically, we expand our analysis to the state-of-the-art Qwen3 model\nobreak\cite{yang2025qwen3technicalreport,qwen}. 
We adopt the widely used C/C++ inference engine llama.cpp\nobreak\cite{llama.cpp} to assess performance in practical scenarios.
This approach allows our evaluation to encompass the entire system from end-to-end (E2E), rather than being limited to single-kernel performance.
This evaluation accounts for host-CPU interaction, data transfer overheads, and complex control flows.

The main contributions of this paper are as follows:
\begin{itemize}
    \item We present the first comprehensive, E2E evaluation of a general-purpose CGRA for modern LLM inference. Our work demonstrates that a non-AI-specialized architecture can achieve superior energy efficiency, improving the Power-Delay Product (PDP) by up to \num{44.4}$\times$ and the Energy-Delay Product (EDP) by up to \num{11.5}$\times$ compared to a high-performance GPU. This finding validates a sustainable hardware approach that avoids the risk of rapid obsolescence associated with task-specific ASICs.
    \item We identify a critical performance bottleneck shift from computation to host-accelerator data transfer. Our system-level analysis using the llama.cpp framework reveals that the decode phase is data-transfer-bound, a finding obscured in kernel-centric evaluations that provides crucial guidance for latency optimization.
    \item We validate CGRAs as a viable architectural solution that balances efficiency and programmability for sustainable LLM inference. Our work establishes that general-purpose reconfigurable hardware can effectively adapt to rapid algorithmic evolution, making it a practical choice for power-constrained deployment.
\end{itemize}



The remainder of this paper is organized as follows. 
Section~\ref{rwork} surveys related work in LLM acceleration and introduces the IMAX architecture.
Section~\ref{proposed} details our implementation methodology, including the execution framework and kernel mapping strategies.
Section~\ref{ex_and_re} presents the experimental results, comparing our accelerator against GPUs on E2E latency and energy-efficiency metrics.
Section~\ref{discussion} analyzes performance bottlenecks by examining on-chip memory impact, execution phase breakdowns, and host system limitations, and discusses directions for future work.
Finally, Section~\ref{conclusions} concludes the paper.

